Die Allgemeine Formel für elektrische Leistung ist:
\begin{equation}
\label{eq:Elektrische_Leistung}
    P = U \cdot I
\end{equation}

\noindent Die am Lastwiderstand anliegende Spannung ergibt sich aus dem Spannungsteiler zu:
\begin{equation}
\label{eq:Spannung_Last}
U = U_0 \frac{R_L}{R_I + R_L}
\end{equation}

\noindent Der durch den Lastwiderstand fließende Strom beträgt:
\begin{equation}
\label{eq:Strom_Last}
I = \frac{U_0}{R_I + R_L}
\end{equation}

\noindent Durch Einsetzen der Gleichungen \ref{eq:Spannung_Last} und \ref{eq:Strom_Last} in Gleichung \ref{eq:Elektrische_Leistung} ergibt sich die im Lastwiderstand umgesetzte Leistung:
\begin{equation}
\label{eq:Leistung_Last_Spannungsquelle}
P_L(R_L) = U_0^2 \frac{R_L}{(R_I + R_L)^2}
\end{equation}